\documentclass[11pt]{article}
\usepackage{amsmath, amssymb, amsthm}
\usepackage[margin=1.2in]{geometry}
\usepackage{hyperref}
\hypersetup{colorlinks=true, linkcolor=blue, urlcolor=blue}
\theoremstyle{definition}
\newtheorem{definition}{Definition}[section]
\newtheorem{theorem}{Theorem}[section]
\newtheorem{lemma}{Lemma}[section]
\newtheorem{remark}{Remark}[section]
\newtheorem{example}{Example}[section]
\newtheorem{corollary}{Corollary}[section]
\begin{document}

\begin{center}
{\Large\textbf{1.1 Topological Preliminaries}}\\
\vspace{0.3cm}
{\normalsize\today}
\end{center}
\vspace{0.5cm}


% --- Page 1 ---
\begin{itemize}
    \item A \textit{topology} $T$ in a set $X$ is a collection of subsets of $X$ (called open sets) satisfying the following axioms:
    \begin{enumerate}
        \item the union of any number of open sets is open
        \item the intersection of any finite number of open sets is open.
        \item $X$, $\emptyset$ are open.
    \end{enumerate}

    A \textit{topological space} is a set $X$ with a topology $T$ in $X$.

    \begin{definition}
        $T = \{\emptyset, X\}$ is \textit{trivial topology}
    \end{definition}

    \begin{definition}
        $T = \mathcal{P}(X)$ is \textit{discrete topology}
    \end{definition}

    \item For $A \subset X$, where $T$ is a topology in $X$. Then $T_A$ is an \textit{induced topology} in $A$
\end{itemize}

% --- Page 2 ---
\[ T_A = \bigcap \{ U \cap A \mid U \in \mathcal{T}_\xi \} \]

\begin{definition}
In a topological spaces, a closed set is a complement of an open set.
\end{definition}

\begin{definition}
For \( E \subseteq X \), \( \overline{E} \) is said to be closure of \( E \), where \( \overline{E} = \bigcap \{ V \mid V \text{ is the set of all closed sets containing } E \} \).
\end{definition}

\begin{definition}
A neighbourhood of \( x \in X \), is any open set containing \( x \).
\end{definition}

\begin{definition}
A mapping \( f \) from \( X_1 \) to \( X_2 \) (\( X_1 \xrightarrow{f} X_2 \)) is an assignment of each \( x \in X_1 \) of an element \( f(x) \in X_2 \) where \( X_1, X_2 \) are topological spaces. (map and mapping are not the same, mapping is a function between two topological spaces).
\end{definition}

% --- Page 3 ---
\begin{itemize}
    \item The mapping $f$ is \underline{continuous} if and only if for every open set $O_2$ in $X_2$, the set $f^{-1}(O_2)$ is an open set in $X_1$. (note $f^{-1}$ need not be a mapping)

    \item For $x \in X_1$, we say that $f$ is continuous at $x$, if to every neighbourhood $(U)$ of $f(x)$, there exists a nbd $V$ of $x$, such that $f(V) \subset U$. ($f$ is continuous at every point of $X_1$)
\end{itemize}

We can show the equivalence of both definitions

\begin{lemma}
    Composition of continuous maps are continuous.
\end{lemma}

\begin{proof}
    If $f: X \rightarrow Y$, $g: Y \rightarrow Z$ are continuous.

    Goal 1: $(g \circ f)^{-1}(V)$ is open in $X$ where $V$ is any open set in $Z$.

    Goal 2: $(g \circ f)^{-1}(V) = f^{-1} \circ g^{-1}(V)$
\end{proof}

% --- Page 4 ---
(i) Every mapping is continuous from $X$ to $Y$ if $X$ has discrete topology. (I showed converse is NOT true).

(ii) If $X$ is a topological space, $Y$ a set with trivial topology, $f: X \rightarrow Y$ is continuous.

* Topological product of $X_1, X_2$ is $X = X_1 \times X_2$ where topology in $X$ is generated by a basis
\[ \mathcal{B} = \{ O_1 \times O_2 \mid O_1 \in \tau_1, O_2 \in \tau_2 \} \]
(Def$^{n}$ of basis add to notes; I'm not writing here)

* If $p_1: X_1 \times X_2 \rightarrow X_1$ and
\[ p_2: X_1 \times X_2 \rightarrow X_2 \] are such that
\[ p_1(x_1 \times x_2) = x_1, \text{ and } p_2(x_1 \times x_2) = x_2 \text{ resp.} \]
Then $p_1, p_2$ are projections of $X_1 \times X_2$.

% --- Page 5 ---
\begin{lemma}
Projections are continuous, open mappings.
\end{lemma}

\begin{proof}
It suffice to show for $p_1$, and for $p_2$ is the same argument.

For open map, observe every open set in $X$ can be written as $U \times V$ where $V$ is an open set in $X_1$, and $V$ is an open set in $X_2$.

Now, $p_1(U \times V) = U$.

For continuity, observe $p_1^{-1}(U) = U \times X_2$ where $U \subset X_1$ is open.
\end{proof}

\begin{remark}
* A family of subsets $\{V_\alpha\}_{\alpha \in I}$ is a \underline{covering} if $\bigcup_{\alpha \in I} V_\alpha = X$. If each $V_\alpha$ is an open subset of $X$, $\{V_\alpha\}_{\alpha \in I}$ is an \underline{open covering} and if $|I| < +\infty$ then $\{V_\alpha\}_{\alpha \in I}$ is a \underline{finite covering}.
\end{remark}

% --- Page 6 ---
\begin{definition}
A topological space $X$ is \underline{compact}, if every open covering of $X$ contains a finite covering.
A subset $A$ is said to be compact, if $A$ is compact in the induced topology.
\end{definition}

\begin{remark}
(i) $\mathbb{R}$ is not compact.

(ii) If $A \subseteq X$, and $|A| < +\infty$, then $A$ is compact.

(iii) Let $a, b \in \mathbb{R}$, $a \leq b$ then $[a, b]$ is compact in $\mathbb{R}$.
\end{remark}

\begin{proof}
Refer to theorem 27.1 and corollary 27.2 of Munkres.
\end{proof}

Let $X$ be a simply ordered set having least upper bound property. Then every closed interval in $X$ is compact. (Check out Prahlad sir's notes on General Topology for alternate path)

% --- Page 7 ---
Given an open covering of \([a, b]\) say, \(\mathcal{A}\).

**STEP 1:** For \(x \in [a, b]\) there exists \(y \in [a, b]\) and \(y > x\) such that \([a, y]\) is covered by at most two elements of \(\mathcal{A}\). (Do case where immediate successor exists and immediate successor doesn't exist).

**STEP 2:**
\[ C = \{ y \in [a, b] \mid [a, y] \text{ are covered by finite elements of } \mathcal{A} \} \]

Observe \(C \neq \emptyset\) by applying STEP 1 for \(x = a\) and apply least upper bound property where \(c = \sup C\).

**STEP 3:** Show \(c \in C\).

**STEP 4:** \(b = C\).

% --- Page 8 ---
* Finite intersection property: In topological space $X$, let $\mathcal{F}_\alpha$ be collection of closed sets. If intersection of any finite subcollection of $\mathcal{F}_\alpha$ is non-empty, then $\bigcap_\alpha F_\alpha \neq \emptyset$.

A space holds FIT iff the space is cpt.

\begin{theorem}
A closed subset of a compact set is compact
\end{theorem}

\begin{proof}
Let $X$ be compact, and $Y$ be a closed subset of $X$.

For any open covering of $Y$, say $\mathcal{A}$, we construct $\mathcal{A} \cup \{X \setminus Y\}$ as an open covering of $X$ as $X \setminus Y$ is open. As $X$ is compact, let $\hat{\mathcal{A}}$ be a finite subcover of $X$. Now, the finite subcover of $Y$ we get from $\hat{\mathcal{A}}$ will not contain $X \setminus Y$, hence finite subcover of $\mathcal{A}$.
\end{proof}

% --- Page 9 ---
\begin{theorem}
A continuous image of a compact set is compact.
\end{theorem}
\begin{proof}
$\checkmark$
\end{proof}

\begin{theorem}
If $X$ is compact and $Y$ is Hausdorff and $f: X \rightarrow Y$ is continuous, then $f(X)$ is closed in $Y$.
\end{theorem}
\begin{proof}
$\checkmark$
\end{proof}

\begin{definition}
A space $X$ is locally compact, if for every $x \in X$, t $\mathbb{R}$
\end{remark}
\begin{proof}
$\mathbb{R}$ is not compact.
\end{proof}

% --- Page 10 ---
For $x \in \mathbb{R}$, there exists $(x-\epsilon, x+\epsilon)$

as a nbhd of $x$ and $\overline{(x-\epsilon, x+\epsilon)} = [x-\epsilon, x+\epsilon]$

is compact.

\begin{definition}
Homeomorphism is a mapping between two topological spaces $X, Y$ say $f$ such that $f$ is bijective map where $f, f^{-1}$ are continuous mappings.
\end{definition}

\textbf{Examples}

Add from C. Chandrashekaran.

% --- Page 11 ---
Note: All topological spaces are \( T_1 \), moving forward unless mentioned otherwise.

\begin{itemize}
    \item Two sets \( M_1 \) and \( M_2 \) are \underline{separated by open sets} if there exist open sets \( O_1 \) and \( O_2 \) containing \( M_1 \) and \( M_2 \) such that \( O_1 \cap O_2 = \emptyset \).

    \item Two sets \( M_1 \) and \( M_2 \) are \underline{separated by a real function} if there exists a continuous real function
    \[ f: X \rightarrow [0, 1] \]
    such that \( f(x) = 0 \) for \( x \in M_1 \) and \( f(x) = 1 \) for \( x \in M_2 \).

    \begin{definition} (\( T_1 \))
        For any disjoint points, there exists an open set containing one and not the other. Equivalently, all singletons are closed.
    \end{definition}

    \begin{definition} (\( T_2 \) (Hausdorff))
        Two distinct points can be separated by open sets.
    \end{definition}

    \begin{definition} (\( T_3 \) (Regular))
        For closed set \( F \) and \( x \notin F \), can be separated by open sets.
    \end{definition}

    \begin{definition} (\( T_4 \) (Normal))
        Disjoint closed sets can be separated
    \end{definition}
\end{itemize}

% --- Page 12 ---
\begin{remark}
$T_4 \Rightarrow T_3 \Rightarrow T_2 \Rightarrow T_1$.
\end{remark}

\textbf{Examples:} From the book.

\begin{lemma}[Urysohn's lemma]
In a normal space, every two disjoint closed sets can be separated by a real function.
\end{lemma}

\begin{theorem}[Tychanoff Theorem]
Product of compact sets are compact.
\end{theorem}

There are two approaches of proving this theorem:
\begin{enumerate}
    \item Zorn's lemma
    \item Transfinite Induction
\end{enumerate}

% --- Page 13 ---
\begin{definition}
Let $X$ be a topological space. A subset $A \subseteq X$ is called
\begin{itemize}
    \item \textbf{open} if $A$ is an open set in $X$,
    \item \textbf{closed} if $X \setminus A$ is open,
    \item \textbf{compact} if every open cover of $A$ has a finite subcover,
    \item \textbf{connected} if there do not exist two non-empty open sets $U, V \subseteq X$ such that $A \subseteq U \cup V$ and $U \cap A \cap V \cap A = \emptyset$.
\end{itemize}
\end{definition}

\begin{theorem}
Let $X$ be a topological space and $A \subseteq X$. Then the following are equivalent:
\begin{enumerate}
    \item $A$ is compact,
    \item Every family of closed subsets of $A$ with the finite intersection property has a non-empty intersection,
    \item Every ultrafilter on $A$ has a non-empty intersection.
\end{enumerate}
\end{theorem}

\begin{lemma}
Let $X$ be a topological space and $A \subseteq X$. If $A$ is compact and $B \subseteq A$ is closed in $A$, then $B$ is compact in $X$.
\end{lemma}

\begin{proof}
Let $\{U_i\}_{i \in I}$ be an open cover of $B$ in $X$. Then $\{U_i \cap A\}_{i \in I}$ is an open cover of $B$ in $A$ (since $B$ is closed in $A$, $A \setminus B$ is open in $A$ and $B$ is closed in $X$). Since $A$ is compact, there exists a finite subcover $\{U_{i_1} \cap A, \ldots, U_{i_n} \cap A\}$. Then $\{U_{i_1}, \ldots, U_{i_n}\}$ is a finite subcover of $B$ in $X$.
\end{proof}

\begin{remark}
The above lemma shows that closed subsets of compact spaces are compact.
\end{remark}

\begin{definition}
A topological space $X$ is called \textbf{Hausdorff} if for every pair of distinct points $x, y \in X$, there exist disjoint open sets $U, V$ such that $x \in U$ and $y \in V$.
\end{definition}

\begin{theorem}
Let $X$ be a compact Hausdorff space. Then every closed subset of $X$ is compact.
\end{theorem}

\begin{proof}
Let $A \subseteq X$ be closed. Then $A$ is compact if and only if every open cover of $A$ has a finite subcover. Let $\{U_i\}_{i \in I}$ be an open cover of $A$. Since $X$ is compact, there exists a finite subcover $\{U_{i_1}, \ldots, U_{i_n}\}$ of $A$ in $X$. Since $A$ is closed, $X \setminus A$ is open, and $\{U_{i_1}, \ldots, U_{i_n}, X \setminus A\}$ is an open cover of $X$. By compactness of $X$, there exists a finite subcover of $X$, which must contain a finite subcover of $A$.
\end{proof}

% illegible

\begin{definition}
A function $f: X \to Y$ between topological spaces is called
\begin{itemize}
    \item \textbf{continuous} if the preimage of every open set in $Y$ is open in $X$,
    \item \textbf{open} if the image of every open set in $X$ is open in $Y$,
    \item \textbf{closed} if the image of every closed set in $X$ is closed in $Y$.
\end{itemize}
\end{definition}

\begin{theorem}
Let $f: X \to Y$ be a continuous function and $A \subseteq X$ be compact. Then $f(A)$ is compact in $Y$.
\end{theorem}

\begin{proof}
Let $\{V_i\}_{i \in I}$ be an open cover of $f(A)$. Then $\{f^{-1}(V_i)\}_{i \in I}$ is an open cover of $A$. Since $A$ is compact, there exists a finite subcover $\{f^{-1}(V_{i_1}), \ldots, f^{-1}(V_{i_n})\}$. Then $\{V_{i_1}, \ldots, V_{i_n}\}$ is a finite subcover of $f(A)$.
\end{proof}

\begin{remark}
This theorem shows that continuous images of compact sets are compact.
\end{remark}

% illegible

\begin{definition}
A topological space $X$ is called \textbf{locally compact} if every point in $X$ has a compact neighborhood.
\end{definition}

% illegible

\begin{theorem}
Let $X$ be a locally compact Hausdorff space. Then every point in $X$ has a neighborhood basis consisting of compact sets.
\end{theorem}

% illegible

\begin{definition}
A topological space $X$ is called \textbf{second-countable} if it has a countable basis for its topology.
\end{definition}

\begin{theorem}
Every second-countable space is separable.
\end{theorem}

\begin{proof}
Let $\{U_n\}_{n \in \mathbb{N}}$ be a countable basis for the topology of $X$. For each $n \in \mathbb{N}$, choose a point $x_n \in U_n$. Then the set $\{x_n\}_{n \in \mathbb{N}}$ is countable and dense in $X$.
\end{proof}

% illegible

\begin{definition}
A topological space $X$ is called \textbf{metrizable} if its topology can be induced by a metric.
\end{definition}

\begin{theorem}
A topological space $X$ is metrizable if and only if it is Hausdorff and second-countable.
\end{theorem}

% illegible

\begin{definition}
A topological space $X$ is called \textbf{paracompact} if every open cover of $X$ has a locally finite open refinement.
\end{definition}

\begin{theorem}
Every paracompact Hausdorff space is normal.
\end{theorem}

% illegible

\begin{definition}
A topological space $X$ is called \textbf{normal} if for every pair of disjoint closed sets $A, B \subseteq X$, there exist disjoint open sets $U, V$ such that $A \subseteq U$ and $B \subseteq V$.
\end{definition}

% illegible

\begin{lemma}
Let $X$ be a normal space and $A \subseteq X$ be closed. Then there exists a continuous function $f: X \to [0, 1]$ such that $f(x) = 0$ for all $x \in A$ and $f(x) = 1$ for all $x \in X \setminus U$, where $U$ is an open set containing $A$.
\end{lemma}

% illegible

\begin{proof}
% illegible
\end{proof}

% illegible

\begin{definition}
A topological space $X$ is called \textbf{completely regular} if for every closed set $A \subseteq X$ and every point $x \notin A$, there exists a continuous function $f: X \to [0, 1]$ such that $f(x) = 0$ and $f(a) = 1$ for all $a \in A$.
\end{definition}

% illegible

\begin{theorem}
Every completely regular space is normal.
\end{theorem}

% illegible

\begin{definition}
A topological space $X$ is called \textbf{completely metrizable} if it is metrizable and complete with respect to some metric inducing its topology.
\end{definition}

% illegible

\begin{theorem}
A topological space $X$ is completely metrizable if and only if it is metrizable and every closed subset of $X$ is complete.
\end{theorem}

% illegible

\begin{definition}
A topological space $X$ is called \textbf{sigma-compact} if it is a countable union of compact sets.
\end{definition}

% illegible

\begin{theorem}
Every sigma-compact space is Lindelöf.
\end{theorem}

% illegible

\begin{definition}
A topological space $X$ is called \textbf{Lindelöf} if every open cover of $X$ has a countable subcover.
\end{definition}

% illegible

\begin{theorem}
Every second-countable space is Lindelöf.
\end{theorem}

% illegible

\begin{definition}
A topological space $X$ is called \textbf{first-countable} if every point in $X$ has a countable neighborhood basis.
\end{definition}

% illegible

\begin{theorem}
Every first-countable space has a countable basis at each point.
\end{theorem}

% illegible

\begin{definition}
A topological space $X$ is called \textbf{hereditarily separable} if every subspace of $X$ is separable.
\end{definition}

% illegible

\begin{theorem}
Every hereditarily separable space is second-countable.
\end{theorem}

% illegible

\begin{definition}
A topological space $X$ is called \textbf{hereditarily Lindelöf} if every subspace of $X$ is Lindelöf.
\end{definition}

% illegible

\begin{theorem}
Every hereditarily Lindelöf space is Lindelöf.
\end{theorem}

% illegible

\begin{definition}
A topological space $X$ is called \textbf{perfect} if every non-empty open set in $X$ contains a non-empty perfect subset.
\end{definition}

% illegible

\begin{theorem}
Every uncountable Polish space contains a perfect subset.
\end{theorem}

% illegible

\begin{definition}
A topological space $X$ is called \textbf{Polish} if it is separable and completely metrizable.
\end{definition}

% illegible

\begin{theorem}
Every closed subset of a Polish space is Polish.
\end{theorem}

% illegible

\begin{definition}
A topological space $X$ is called \textbf{analytic} if it is a continuous image of a Polish space.
\end{definition}

% illegible

\begin{theorem}
Every analytic space is separable.
\end{theorem}

% illegible

\begin{definition}
A topological space $X$ is called \textbf{co-analytic} if its complement in some Polish space is analytic.
\end{definition}

% illegible

\begin{theorem}
Every co-analytic subset of a Polish space is either countable or contains a perfect subset.
\end{theorem}

% illegible

\begin{definition}
A topological space $X$ is called \textbf{universal} for a class of spaces if every space in that class is homeomorphic to a subspace of $X$.
\end{definition}

% illegible

\begin{theorem}
There exists a universal Polish space for all separable metrizable spaces.
\end{theorem}

% illegible

\begin{definition}
A topological space $X$ is called \textbf{countably compact} if every countable open cover of $X$ has a finite subcover.
\end{definition}

% illegible

\begin{theorem}
Every compact space is countably compact.
\end{theorem}

% illegible

\begin{definition}
A topological space $X$ is called \textbf{Lindelöf} if every open cover of $X$ has a countable subcover.
\end{definition}

% illegible

\begin{theorem}
Every Lindelöf space is countably compact.
\end{theorem}

% illegible

\begin{definition}
A topological space $X$ is called \textbf{pseudocompact} if every continuous real-valued function on $X$ is bounded.
\end{definition}

% illegible

\begin{theorem}
Every pseudocompact space is countably compact.
\end{theorem}

% illegible

\begin{definition}
A topological space $X$ is called \textbf{sequentially compact} if every sequence in $X$ has a convergent subsequence.
\end{definition}

% illegible

\begin{theorem}
Every compact metrizable space is sequentially compact.
\end{theorem}

% illegible

\begin{definition}
A topological space $X$ is called \textbf{sequentially closed} if it contains all limits of sequences of its points.
\end{definition}

% illegible

\begin{theorem}
Every sequentially compact space is countably compact.
\end{theorem}

% illegible

\begin{definition}
A topological space $X$ is called \textbf{sequentially Lindelöf} if every sequence in $X$ has a countable closure.
\end{definition}

% illegible

\begin{theorem}
Every sequentially Lindelöf space is Lindelöf.
\end{theorem}

% illegible

\begin{definition}
A topological space $X$ is called \textbf{hereditarily separable} if every subspace of $X$ is separable.
\end{definition}

% illegible

\begin{theorem}
Every hereditarily separable space is second-countable.
\end{theorem}

% illegible

\begin{definition}
A topological space $X$ is called \textbf{hereditarily Lindelöf} if every subspace of $X$ is Lindelöf.
\end{definition}

% illegible

\begin{theorem}
Every hereditarily Lindelöf space is Lindelöf.
\end{theorem}

% illegible

\begin{definition}
A topological space $X$ is called \textbf{hereditarily normal} if every subspace of $X$ is normal.
\end{definition}

% illegible

\begin{theorem}
Every hereditarily normal space is completely normal.
\end{theorem}

% illegible

\begin{definition}
A topological space $X$ is called \textbf{completely normal} if for every pair of disjoint closed sets $A, B \subseteq X$, there exists a continuous function $f: X \to [0, 1]$ such that $f(A) = \{0\}$ and $f(B) = \{1\}$.
\end{definition}

% illegible

\begin{theorem}
Every completely normal space is normal.
\end{theorem}

% illegible

\begin{definition}
A topological space $X$ is called \textbf{hereditarily paracompact} if every subspace of $X$ is paracompact.
\end{definition}

% illegible

\begin{theorem}
Every hereditarily paracompact space is paracompact.
\end{theorem}

% illegible

\begin{definition}
A topological space $X$ is called \textbf{hereditarily metrizable} if every subspace of $X$ is metrizable.
\end{definition}

% illegible

\begin{theorem}
Every hereditarily metrizable space is metrizable.
\end{theorem}

% illegible

\begin{definition}
A topological space $X$ is called \textbf{hereditarily second-countable} if every subspace of $X$ is second-countable.
\end{definition}

% illegible

\begin{theorem}
Every hereditarily second-countable space is second-countable.
\end{theorem}

% illegible

\begin{definition}
A topological space $X$ is called \textbf{hereditarily separable} if every subspace of $X$ is separable.
\end{definition}

% illegible

\begin{theorem}
Every hereditarily separable space is second-countable.
\end{theorem}

% illegible

\begin{definition}
A topological space $X$ is called \textbf{hereditarily Lindelöf} if every subspace of $X$ is Lindelöf.
\end{definition}

% illegible

\begin{theorem}
Every hereditarily Lindelöf space is Lindelöf.
\end{theorem}

% illegible

\begin{definition}
A topological space $X$ is called \textbf{hereditarily compact} if every subspace of $X$ is compact.
\end{definition}

% illegible

\begin{theorem}
Every hereditarily compact space is compact.
\end{theorem}

% illegible

\begin{definition}
A topological space $X$ is called \textbf{hereditarily connected} if every subspace of $X$ is connected.
\end{definition}

% illegible

\begin{theorem}
Every hereditarily connected space is connected.
\end{theorem}

% illegible

\begin{definition}
A topological space $X$ is called \textbf{hereditarily path-connected} if every subspace of $X$ is path-connected.
\end{definition}

% illegible

\begin{theorem}
Every hereditarily path-connected space is path-connected.
\end{theorem}

% illegible

\begin{definition}
A topological space $X$ is called \textbf{hereditarily simply connected} if every subspace of $X$ is simply connected.
\end{definition}

% illegible

\begin{theorem}
Every hereditarily simply connected space is simply connected.
\end{theorem}

% illegible

\begin{definition}
A topological space $X$ is called \textbf{hereditarily locally compact} if every subspace of $X$ is locally compact.
\end{definition}

% illegible

\begin{theorem}
Every hereditarily locally compact space is locally compact.
\end{theorem}

% illegible

\begin{definition}
A topological space $X$ is called \textbf{hereditarily locally connected} if every subspace of $X$ is locally connected.
\end{definition}

% illegible

\begin{theorem}
Every hereditarily locally connected space is locally connected.
\end{theorem}

% illegible

\begin{definition}
A topological space $X$ is called \textbf{hereditarily locally path-connected} if every subspace of $X$ is locally path-connected.
\end{definition}

% illegible

\begin{theorem}
Every hereditarily locally path-connected space is locally path-connected.
\end{theorem}

% illegible

\begin{definition}
A topological space $X$ is called \textbf{hereditarily locally simply connected} if every subspace of $X$ is locally simply connected.
\end{definition}

% illegible

\begin{theorem}
Every hereditarily locally simply connected space is locally simply connected.
\end{theorem}

% illegible

\begin{definition}
A topological space $X$ is called \textbf{hereditarily first-countable} if every subspace of $X$ is first-countable.
\end{definition}

% illegible

\begin{theorem}
Every hereditarily first-countable space is first-countable.
\end{theorem}

% illegible

\begin{definition}
A topological space $X$ is called \textbf{hereditarily second-countable} if every subspace of $X$ is second-countable.
\end{definition}

% illegible

\begin{theorem}
Every hereditarily second-countable space is second-countable.
\end{theorem}

% illegible

\begin{definition}
A topological space $X$ is called \textbf{hereditarily separable} if every subspace of $X$ is separable.
\end{definition}

% illegible

\begin{theorem}
Every hereditarily separable space is second-countable.
\end{theorem}

% illegible

\begin{definition}
A topological space $X$ is called \textbf{hereditarily Lindelöf} if every subspace of $X$ is Lindelöf.
\end{definition}

% illegible

\begin{theorem}
Every hereditarily Lindelöf space is Lindelöf.
\end{theorem}

% illegible

\begin{definition}
A topological space $X$ is called \textbf{hereditarily normal} if every subspace of $X$ is normal.
\end{definition}

% illegible

\begin{theorem}
Every hereditarily normal space is completely normal.
\end{theorem}

% illegible

\begin{definition}
A topological space $X$ is called \textbf{hereditarily completely normal} if every subspace of $X$ is completely normal.
\end{definition}

% illegible

\begin{theorem}
Every hereditarily completely normal space is completely normal.
\end{theorem}

% illegible

\begin{definition}
A topological space $X$ is called \textbf{hereditarily paracompact} if every subspace of $X$ is paracompact.
\end{definition}

% illegible

\begin{theorem}
Every hereditarily paracompact space is paracompact.
\end{theorem}

% illegible

\begin{definition}
A topological space $X$ is called \textbf{hereditarily metrizable} if every subspace of $X$ is metrizable.
\end{definition}

% illegible

\begin{theorem}
Every hereditarily metrizable space is metrizable.
\end{theorem}

% illegible

\begin{definition}
A topological space $X$ is called \textbf{hereditarily second-countable} if every subspace of $X$ is second-countable.
\end{definition}

% illegible

\begin{theorem}
Every hereditarily second-countable space is second-countable.
\end{theorem}

% illegible

\begin{definition}
A topological space $X$ is called \textbf{hereditarily separable} if every subspace of $X$ is separable.
\end{definition}

% illegible

\begin{theorem}
Every hereditarily separable space is second-countable.
\end{theorem}

% illegible

\begin{definition}
A topological space $X$ is called \textbf{hereditarily Lindelöf} if every subspace of $X$ is Lindelöf.
\end{definition}

% illegible

\begin{theorem}
Every hereditarily Lindelöf space is Lindelöf.
\end{theorem}

% illegible

\begin{definition}
A topological space $X$ is called \textbf{hereditarily compact} if every subspace of $X$ is compact.
\end{definition}

% illegible

\begin{theorem}
Every hereditarily compact space is compact.
\end{theorem}

% illegible

\begin{definition}
A topological space $X$ is called \textbf{hereditarily connected} if every subspace of $X$ is connected.
\end{definition}

% illegible

\begin{theorem}
Every hereditarily connected space is connected.
\end{theorem}

% illegible

\begin{definition}
A topological space $X$ is called \textbf{hereditarily path-connected} if every subspace of $X$ is path-connected.
\end{definition}

% illegible

\begin{theorem}
Every hereditarily path-connected space is path-connected.
\end{theorem}

% illegible

\begin{definition}
A topological space $X$ is called \textbf{hereditarily simply connected} if every subspace of $X$ is simply connected.
\end{definition}

% illegible

\begin{theorem}
Every hereditarily simply connected space is simply connected.
\end{theorem}

% illegible

\begin{definition}
A topological space $X$ is called \textbf{hereditarily locally compact} if every subspace of $X$ is locally compact.
\end{definition}

% illegible

\begin{theorem}
Every hereditarily locally compact space is locally compact.
\end{theorem}

% illegible

\begin{definition}
A topological space $X$ is called \textbf{hereditarily locally connected} if every subspace of $X$ is locally connected.
\end{definition}

% illegible

\begin{theorem}
Every hereditarily locally connected space is locally connected.
\end{theorem}

% illegible

\begin{definition}
A topological space $X$ is called \textbf{hereditarily locally path-connected} if every subspace of $X$ is locally path-connected.
\end{definition}

% illegible

\begin{theorem}
Every hereditarily locally path-connected space is locally path-connected.
\end{theorem}

% illegible

\begin{definition}
A topological space $X$ is called \textbf{hereditarily locally simply connected} if every subspace of $X$ is locally simply connected.
\end{definition}

% illegible

\begin{theorem}
Every hereditarily locally simply connected space is locally simply connected.
\end{theorem}

% illegible

\begin{definition}
A topological space $X$ is called \textbf{hereditarily first-countable} if every subspace of $X$ is first-countable.
\end{definition}

% illegible

\begin{theorem}
Every hereditarily first-countable space is first-countable.
\end{theorem}

% illegible

\begin{definition}
A topological space $X$ is called \textbf{hereditarily second-countable} if every subspace of $X$ is second-countable.
\end{definition}

% illegible

\begin{theorem}
Every hereditarily second-countable space is second-countable.
\end{theorem}

% illegible

\begin{definition}
A topological space $X$ is called \textbf{hereditarily separable} if every subspace of $X$ is separable.
\end{definition}

% illegible

\begin{theorem}
Every hereditarily separable space is second-countable.
\end{theorem}

% illegible

\begin{definition}
A topological space $X$ is called \textbf{hereditarily Lindelöf} if every subspace of $X$ is Lindelöf.
\end{definition}

% illegible

\begin{theorem}
Every hereditarily Lindelöf space is Lindelöf.
\end{theorem}

% illegible

\begin{definition}
A topological space $X$ is called \textbf{hereditarily normal} if every subspace of $X$ is normal.
\end{definition}

% illegible

\begin{theorem}
Every hereditarily normal space is completely normal.
\end{theorem}

% illegible

\begin{definition}
A topological space $X$ is called \textbf{hereditarily completely normal} if every subspace of $X$ is completely normal.
\end{definition}

% illegible

\begin{theorem}
Every hereditarily completely normal space is completely normal.
\end{theorem}

% illegible

\begin{definition}
A topological space $X$ is called \textbf{hereditarily paracompact} if every subspace of $X$ is paracompact.
\end{definition}

% illegible

\begin{theorem}
Every hereditarily paracompact space is paracompact.
\end{theorem}

% illegible

\begin{definition}
A topological space $X$ is called \textbf{hereditarily metrizable} if every subspace of $X$ is metrizable.
\end{definition}

% illegible

\begin{theorem}
Every hereditarily metrizable space is metrizable.
\end{theorem}

\end{document}